% SOURCE_AUTHORITY: sga5_fr_workpass.tex, lines 13141--13161 (LF-normalized unit).
% SOURCE_SHA256: 791F4EFFC5E02832D5D77ED03518C8156D6F07E4C8238B03545DB93D883FBB28
Seja, para isso, $f$ um gerador topológico de $\bar\pi$. Pode-se supor que as componentes $M^i$ de $M$ são $\bar\pi$-módulos tais que
\[
H^j(\bar\pi,M^i)=0\quad\text{se }j\ge1.
\]
Existe, então, um triângulo distinguido
\[
\begin{tikzcd}[column sep=huge,row sep=large]
& M^\bullet \arrow[dl] &\\
\R\Gamma^{\bar\pi}(M^\bullet) \arrow[rr] && M^\bullet \arrow[ul,"1-f"]
\end{tikzcd}
\tag{7.15}
\]
Com efeito, por um lado, é claro que $\R\Gamma^{\bar\pi}=\Ker(1-f)$. Por outro lado, $1-f:M^\bullet\to M^\bullet$ é sobrejetivo. Com efeito,
\[
H^1(\bar\pi,M)=M^\bullet/(1-f)M^\bullet
\]
(veja-se, para este cálculo, [4] p. 197, em um caso sensivelmente análogo), e resulta, portanto, $M^\bullet=(1-f)M^\bullet$. O triângulo (7.15) fornece simultaneamente $\R\Gamma^{\bar\pi}(M^\bullet)\in\Ob(\Delta)$ e
\[
\clK{\Delta}(\R\Gamma^{\bar\pi}(M^\bullet))=0,
\]
que demonstra 7.14.
\end{prooftext}
